% Generated draft: requires mathematical and source review.
\documentclass[11pt]{article}
\usepackage[margin=25mm]{geometry}
\usepackage{fontspec}
\setmainfont{Latin Modern Roman}
\usepackage{amsmath,amssymb,mathtools,unicode-math}
\setmathfont{Latin Modern Math}
\usepackage{enumitem,needspace,xurl,hyperref}
\setlength{\parindent}{0pt}
\setlength{\parskip}{5pt}
\setlength{\emergencystretch}{3em}
\newcommand{\N}{\mathbb{N}}
\newcommand{\Z}{\mathbb{Z}}
\newcommand{\Q}{\mathbb{Q}}
\newcommand{\R}{\mathbb{R}}
\newcommand{\C}{\mathbb{C}}
\newcommand{\F}{\mathbb{F}}
\newcommand{\A}{\mathbb{A}}
\newcommand{\PP}{\mathbb{P}}
\newcommand{\E}{\mathbb{E}}
\newcommand{\eps}{\varepsilon}
\newcommand{\dd}{\,\mathrm{d}}
\begin{document}
\section*{UnsolvedMath ID: 11100017}
\textbf{Source problem code:} catalogue entry
\subsection{Definitions and mathematical statement}
Let \(E\) be the ring spectrum and \(S\) the group in the source's Morava-theory setup. Suppose the intended action of \(S\) on \(E\) is by ring-spectrum automorphisms. An \(E\)--\(S\) module should consist of an \(E\)-module spectrum \(M\) and a compatible \(S\)-action. The compatibility is semilinear:
\[
 s(e\cdot m)=s(e)\cdot s(m),\qquad s\in S,
\]
interpreted as commutativity of the action diagrams with the required homotopy coherence. If \(S\) acts trivially on \(E\), this specializes to commuting actions, a different special case. Morphisms should be \(E\)-linear and \(S\)-equivariant in the chosen homotopical sense.

\emph{Question (Mark Hovey).} Can one form a suitable category of such \(E\)--\(S\) module spectra, and is that framework worthwhile? A complete answer must select the spectrum and module model, ring structure, group and topology, continuity and coherence requirements, morphisms and weak equivalences, and demonstrate the intended utility. Merely forming a category of strict algebraic actions does not by itself fix these homotopical issues. The source expressly cautions that formal choices can consume substantial effort.

The supplied summary interprets \(E\) as height-\(n\) Morava \(E_n\) and \(S\) as a stabilizer-type group, and marks the construction solved through modern equivariant-module/homotopy-fixed-point frameworks. The exact prime, height, ordinary versus extended stabilizer, continuous/profinite action model, and \(K(n)\)-localization convention are not supplied. An \(E_n\)-module category must not be identified with the entire \(K(n)\)-local stable category, nor can a naive discrete action silently replace a continuous profinite action. The summary's homotopy-fixed-point and higher-real-\(K\)-theory context is retained without claiming it establishes every version of the original categorical and usefulness question.
\subsection{Short English statement}
Construct a homotopically appropriate category of module spectra with compatible, generally semilinear group actions, and explain its use in the intended Morava setting. The registry labels this solved, but the ring, stabilizer topology, continuity, and localization conventions require a precise match.
\subsection{Sources}
\begin{itemize}
\item Local UnsolvedMath registry, permanent ID 11100017, source code catalogue entry; proposer Mark Hovey, ``Morava K- and E-theory,'' item 7.
\item \emph{Hovey's Algebraic Topology Problem List}, surviving mirror: \url{https://www-users.cse.umn.edu/~tlawson/hovey/}. Recorded access 2026-07-29; extraction html-ordered-list. The mirror contains 65 top-level items across nine category pages.
\item The supplied summary names Devinatz--Hopkins homotopy fixed points, higher real \(K\)-theories \(EO_n\), and Lurie's \emph{Higher Algebra}/infinity-category framework. It provides no exact theorem or full bibliography and ends with truncated text.
\item Supplied status solved and SOLVED-IN-LITERATURE. Its summary states that compatible-action \(E_n\)-module categories exist, but includes an overbroad comparison with all \(K(n)\)-local spectra and incomplete formal specifications. Generic background still requires a current primary-literature audit and review of redistribution terms, both NEEDS\_REVIEW.
\end{itemize}
\end{document}
